\printedpage{44}
% head: A. Grothendieck : Dérivateurs, Ch. II, Édition des Universités de Montpellier II et Paris VII

We shall put, for a subset $A$ of $\cP$ resp. $B$ of $\cQ$
\[
A^{\natural} = (A^{\cup})^{\cap}, \qquad B^{\flat} = (B^{\cap})^{\cup}
\tag*{(5.4)}
\]
so that we shall have
\[
A \subseteq A^{\natural}, \qquad B \subseteq B^{\flat}
\tag*{(5.5)}
\]
We shall say that $A$ is $\cR$-\textit{closed} if $A = A^{\natural}$ (or again if $A$ is of the form $B'^{\cap}$ for $B' \subseteq \cQ$, cf. 5.6), that
\msp{71}
$B$ is $\cR$-\textit{closed} if $B = B^{\flat}$. We shall have for every $A \subseteq \cP$, $B \subseteq \cQ$
\[
\left\{
\begin{array}{rcll}
(A^{\cup})^{\flat} & = & A^{\cup}, & \text{i.e. } A^{\cup} \text{ is } \cR\text{-closed,}\\[1ex]
(B^{\cap})^{\natural} & = & B^{\cap}, & \text{i.e. } B^{\cap} \text{ is } \cR\text{-closed.}
\end{array}
\right.
\tag*{(5.6)}
\]
In particular
\[
(A^{\natural})^{\natural} = A^{\natural}, \qquad (B^{\flat})^{\flat} = B^{\flat},
\tag*{(5.7)}
\]
i.e. the $\cR$-closure $A^{\natural}$ of $A \subseteq \cP$ is $\cR$-closed, and likewise the $\cR$-closure $B^{\flat}$ of $B \subseteq \cQ$. % [sic] cloture -> clôture, accent circonflexe manquant (cf. « clôture » dans la même ligne)

In fact, I realize that I am mixing two different kinds of facts. The relations (5.2), and (5.4) to (5.7), do not depend on the datum of order relations on $\cP$, $\cQ$, nor on axioms connecting these with the given relation $\cR$, but only on this relation, which need not be subject to any condition. Let us note that % [sic] la donné -> la donnée, accord manquant
\[
\emptyset^{\cup} = \cQ, \qquad \emptyset^{\cap} = \cP,
\tag*{(5.8)}
\]
whereas
\[
\emptyset^{\natural} = \cQ^{\cap}, \qquad \emptyset^{\flat} = \cP^{\cup}
\tag*{(5.9)}
\]
are respectively the smallest $\cR$-closed subset of $\cP$ and of $\cQ$.

The map
\[
A \longmapsto A^{\natural}
\tag*{(5.10)}
\]
is an \textit{increasing} (and \textit{idempotent}) map from $\fP(\cP)$ to itself, thus having the property
\msp{72}
\[
(\bigcup_i A_i)^{\natural} = (\bigcup_i A_i^{\natural})^{\natural}.
\]

\printedpage{45}
% head: A. Grothendieck : Dérivateurs, Ch. II, Édition des Universités de Montpellier II et Paris VII

For every $A \subseteq \cP$, $A^{\natural}$ is the smallest $\cR$-closed subset of $\cP$ containing $A$, and likewise for every $B \subseteq \cQ$, $B^{\flat}$ is the smallest $\cR$-closed subset of $Q$ containing $B$. We have moreover (this should even have been said before saying that $A \longmapsto A^{\natural}$ is increasing) % [sic] de $Q$ -> de $\mathcal{Q}$, lettre droite au lieu de la lettre calligraphique ; [sic] contentant -> contenant, faute de frappe
\[
\left\{
\begin{array}{rcl}
A \subseteq A' & \Longrightarrow & A'^{\cup} \subseteq A^{\cup}\\[1ex]
B \subseteq B' & \Longrightarrow & B'^{\cap} \subseteq B^{\cap},
\end{array}
\right.
\tag*{(5.11)}
\]
i.e. the maps
\[
\begin{tikzcd}
\fP(\cP) \arrow[r, shift left, "\beta"] & \fP(\cQ) \arrow[l, shift left, "\alpha"]
\end{tikzcd}
\qquad
\begin{array}{l}
A \longmapsto A^{\cup}\\
B \longmapsto B^{\cap}
\end{array}
\tag*{(5.12)}
\]
are decreasing (which implies that the two composites $\alpha\beta = (A \longmapsto A^{\natural})$ and $\beta\alpha = (B \longmapsto B^{\flat})$ are increasing). Their restrictions to the subsets $\fP_{\cR}(\cP)$ and $\fP_{\sR}(\cQ)$ of the $\cR$-closed subsets of $\cP$ resp. of $\cQ$, are, by definition of these as
\[
\begin{array}{rcl}
\fP_{\cR}(\cP) & = & \{A \subseteq \cP \mid A^{\natural} \ (\overset{\text{def}}{=} \alpha\beta(A)) = A\}\\
\fP_{\sR}(\cQ) & = & \{B \subseteq \cQ \mid B^{\flat} \ (\overset{\text{def}}{=} \beta\alpha(B)) = B\},
\end{array}
\]
inverse to one another. We thus find

\textbf{Proposition 5.1.} \textit{Let $\cP$, $\cQ$ be two sets, $\cR \subseteq \cP \times \cQ$ a relation. Then the maps $\alpha$, $\beta$ of (5.12) induce bijections inverse to one}
\msp{73}
\textit{another, between the set $\mathit{Cl}_{\cR}(\cP)$ of the $\cR$-closed subsets of $\cP$, and the set $\mathit{Cl}_{\sR}(\cQ)$ of the ${}^{s}\cR$-closed subsets (we shall say $\cR$-closed by abuse of language) of $\cQ$. These bijections are decreasing, hence antiisomorphisms for the natural order relations (of inclusion).}

\textbf{Corollary 5.2.} \textit{Let $(A_i)_i$ be a family of $\cR$-closed subsets of $\cP$. Then $\bigcap A_i$ is a $\cR$-closed subset of $\cP$, and}
\[
(\bigcap_i A_i)^{\cup} = \underbrace{(\bigcup_i A_i^{\cup})^{\flat}}_{\substack{\text{sup of the}\\ \text{subsets } \cR\text{-closed}\\ A_i^{\cup} \text{ of } \cQ}} \quad .
\tag*{(5.2.1)}
\]

But the relation $\cR$ ``extends'' to a relation $\bar{\cR}$ between
\[
\bar{\cP} = \fP(\cP), \qquad \bar{\cQ} = \fP(\cQ)
\tag*{(5.13)}
\]
by the prescription
\[
\underbrace{\bar{\cR}}_{\subseteq \bar{\cP} \times \bar{\cQ}} = \{(A,B) \in \bar{\cP} \times \bar{\cQ} \mid B \subseteq A^{\cup},\ \text{i.e. } A \subseteq B^{\cap},\ \text{i.e. } A \times B \subseteq \cR\}.
\tag*{(5.14)}
\]
The sets $\bar{\cP}$, $\bar{\cQ}$ are stable under arbitrary sups (union of subsets of $\cP$ resp. $\cQ$), and the relation $\bar{\cR}$ on these ordered sets satisfies condition (D$_0$)
\printedpage{46}
% head: A. Grothendieck : Dérivateurs, Ch. II, Édition des Universités de Montpellier II et Paris VII
\msp{74}
above (p. 69), in an obvious way: for fixed $A$, the $B$ such that $(A,B) \in \bar{\cR}$ are those which are contained in $A^{\cup}$, so the set $\bar{\cR}(A)$ is nothing other than $\bar{\cQ}/A^{\cup}$, and dually ${}^{s}\bar{\cR}(B)$ is nothing other than $\bar{\cP}/B^{\cap}$. Thus
\[
\beta(A) = A^{\cup}, \qquad \alpha(B) = B^{\cap}.
\tag*{(5.14)} % [sic] (5.14) -> numéro distinct, numéro déjà employé pour la formule (5.14) de la page précédente
\]
If I thought of introducing axiom (D$_0$), for a pair of arbitrary ordered sets, it is because of the cases where these are \textit{not} of the form $\fP(\cP)$, $\fP(\cQ)$, and (even if they were) where the basic relation is not deduced as above from an arbitrary relation between sets $\cP$, $\cQ$.

I therefore return to the axiomatics of the beginning, with $\cP$, $\cQ$ ordered, and $\cR$ satisfying (D). In all the cases I can think of, $\cP$ and $\cQ$ are stable under arbitrary sups, so (D) is equivalent to (D$_0$). Let us note that the map
\[
\cP \longrightarrow \fP(\cP), \qquad p \longmapsto \cP/p
\]
is injective, and likewise for $\cQ \longrightarrow \fP(\cQ)$. Thus, knowledge of $A \longmapsto A^{\cup}$ is equivalent to that of % [sic] connaisance -> connaissance, faute d'orthographe
\[
\begin{array}{rcl}
A & \longmapsto & \beta(A)\\[0.5ex]
\beta : \fP(\cP) & \longrightarrow & \cQ,
\end{array}
\]
and likewise that of $B \longmapsto B^{\cap}$ is equivalent to
\msp{75}
that of
\[
\alpha : \fP(\cQ) \longrightarrow \cP.
\]
In fact, it suffices to give oneself
\[
\begin{array}{rcccl}
\beta_0 : & \cP & \longrightarrow & \cQ & \qquad \text{maps}\\
\alpha_0 : & \cQ & \longrightarrow & \cP & \qquad \textit{decreasing}
\end{array}
\tag*{(5.15)}
\]
defined by
\[
\left\{
\begin{array}{rcl}
\beta_0(p) & = & \beta(\{p\})\\[1ex]
\alpha_0(q) & = & \alpha(\{q\}),
\end{array}
\right.
\]
since we deduce $\beta$, $\alpha$ from them by
\[
\left\{
\begin{array}{rcl}
\beta(A) & = & \inf\nolimits_{p \in A} \beta_0(p)\\[1ex]
\alpha(B) & = & \inf\nolimits_{q \in B} \alpha_0(q).
\end{array}
\right.
\tag*{(5.16)}
\]
The relation $\cR$ is deduced from the datum of $\beta_0$, or at one's choice of $\alpha_0$, by
\[
(p,q) \in \cR \iff q \leq \beta_0(p) \iff p \leq \alpha_0(q).
\tag*{(5.17)}
\]

\printedpage{47}
% head: A. Grothendieck : Dérivateurs, Ch. II, Édition des Universités de Montpellier II et Paris VII

This therefore implies that $\alpha_0$, $\beta_0$ determine one another, by the formulas
\[
\left\{
\begin{array}{rcl}
\alpha_0(q) & = & \sup \underbrace{\{p \in \cP \mid q \leq \beta_0(p)\}}_{\{q\}^{\cap}}\\[3ex]
\beta_0(p) & = & \sup \underbrace{\{q \in \cQ \mid p \leq \alpha_0(q)\}}_{\{p\}^{\cup}}.
\end{array}
\right.
\tag*{(5.18)}
\]
It remains to ask which maps $\alpha_0 : \cP \longrightarrow \cQ$, or $\beta_0 : \cQ \longrightarrow \cP$, correspond to a \textit{``dualizing relation''} % [sic] \alpha_0 : \cP \to \cQ, \beta_0 : \cQ \to \cP -> \beta_0 : \cP \to \cQ, \alpha_0 : \cQ \to \cP, les rôles de \alpha_0 et \beta_0 sont échangés par rapport à (5.15)
\msp{76}
$\cR$ on $\cP \times \cQ$. Here are the conditions on $\beta_0 : \cP \longrightarrow \cQ$:
\begin{itemize}[leftmargin=4em, labelsep=0.8em]
\item[\circled{a}] $\beta_0$ is \textit{decreasing},
\end{itemize}
which implies that the relation
\[
\underbrace{\cR_{\beta_0}}_{\subseteq \cP \times \cQ} = \{(p,q) \in \cP \times \cQ \mid q \leq \beta_0(p)\}
\]
satisfies condition (D$_0$) on the left side, and we still have to check it on the right side, i.e. look, for fixed $q$, at the set
\[
\{q\}^{\cap} = \{p \in \cP \mid q \leq \beta_0(p)\}.
\]
As $\beta_0$ is decreasing, we see that $p \in \{q\}^{\cap}$, and $p' \leq p \Longrightarrow p' \in \{q\}^{\cap}$, i.e. $\{q\}^{\cap}$ is \textit{open.} We must still check (assuming the existence of arbitrary sups in $\cP$) that this set is stable under arbitrary sups - it will follow that it is of the form $\cP/\alpha(q)$. This is equivalent to the condition
\begin{itemize}[leftmargin=4em, labelsep=0.8em]
\item[\circled{b}] $\beta_0$ transforms arbitrary sups in $\cP$ into infs in $\cQ$.
\end{itemize}
If we regard $\beta_0$ as a contravariant functor $\cP^{o} \longrightarrow \cQ$, this means that it commutes with inverse limits, or again that it has a left adjoint $\cQ \longrightarrow \cP^{o}$, which commutes with direct limits, and corresponds
\msp{77}
to the decreasing map $\alpha_0 : \cQ \longrightarrow \cP$, transforming sups into infs.

Here it is convenient to consider $\natural$ and $\flat$ as operations on $\cP$ resp. $\cQ$, by
\[
p^{\natural} = \alpha_0\beta_0(p), \qquad q^{\flat} = \beta_0\alpha_0(q).
\tag*{(5.19)}
\]
The relation between these endomorphisms, and the operation of the same name on $\fP(\cP)$, $\fP(\cQ)$, is given by
\[
\left\{
\begin{array}{rcl}
p^{\natural} & = & \text{Greatest element of } \{p\}^{\natural}\\[1ex]
A^{\natural} & = & \sup\nolimits_{p \in A} p^{\natural},
\end{array}
\right.
\tag*{(5.20)}
\]
\printedpage{48}
% head: A. Grothendieck : Dérivateurs, Ch. II, Édition des Universités de Montpellier II et Paris VII
and symmetrically on the side of $\cQ$.

We say that $p$ is $\cR$-\textit{closed,} or $\beta_0$-closed, or $\alpha_0$-closed, if $p = p^{\natural}$, and likewise $q \in \cQ$ is said to be $\cR$-\textit{closed} etc. if $q = q^{\flat}$. The maps
\[
\left\{
\begin{array}{ccccccc}
p & \longmapsto & p^{\natural} & : & \cP & \longrightarrow & \cP\\[1ex]
q & \longmapsto & q^{\flat} & : & \cQ & \longrightarrow & \cQ
\end{array}
\right.
\tag*{(5.21)}
\]
are \textit{increasing} and \textit{idempotent.} We can characterize $p^{\natural}$ as the smallest closed element of $\cP$ lying above $p$, and likewise for $q^{\flat}$. We then find, by (5.19), that $\alpha_0$ and $\beta_0$ induce bijections inverse to one another,
\msp{78}
between the set of closed elements of $\cP$, and the set of closed elements of $\cQ$
\[
\begin{tikzcd}
\cP^{\natural} \arrow[r, shift left, "\bar{\beta}_0"] \arrow[r, phantom, "\scriptstyle\sim"] & \cQ^{\flat}. \arrow[l, shift left, "\bar{\alpha}_0"]
\end{tikzcd}
\tag*{(5.22)}
\]
But in practice, the maps $\alpha_0$, $\beta_0$ are not obtained directly, but we have at our disposal the relation $\cR$, satisfying condition (D$_0$), from which we deduce $\alpha_0$, $\beta_0$. In all the cases I am thinking of here, either $\cP$, $\cQ$ are of the form $\fP(P_0)$, $\fP(Q_0)$, and $\cR$ is deduced from a relation $R_0 \subseteq P_0 \times Q_0$ as above (by far the most frequent, it seems to me), or else $\cP$ and $\cQ$ appear as suitable subsets of sets $\fP(P_0)$, $\fP(Q_0)$, equipped with the induced order, and the relation $\cR$ on $\cP \times \cQ$ is obtained in a more complex way.

\noindent\textbf{Examples.}
\begin{enumerate}[label={\arabic*)}]
\item Relations of \textit{orthogonality} between subsets of two modules over a ring, equipped with a duality pairing. In the case of vector \texttt{[spaces]} of finite dimension and of a true duality, the closed subsets are the vector subspaces, in the infinite case, they are the vector subspaces closed for the ``weak topology'' defined by the pairing.
\item Relations of \textit{orthogonality} between subsets of two locally compact commutative topological groups dual to one another, in the sense of Pontrjagin.
\msp{79}

The closed subsets are the closed subgroups.
\item For two vector spaces ``in duality'' over a complete valued field, such as $\mathbf{R}$ or $\mathbf{C}$, the relation % [sic] où -> ou, accent grave fautif
\[
|\langle x, x' \rangle| \leq 1.
\]
In the real or complex case, the closed subsets are the convex circled [\,?] weakly closed subsets.
\item For real vector spaces, the relation of semi-polarity
\[
\langle x, x' \rangle \geq -1.
\]
The closed subsets are the convex weakly closed subsets containing the origin.
\end{enumerate}